\subsection{Technical Competency Achievement}

This project has served as a practical means for the development and consolidation of several key technical competencies within the field of Software Engineering. This section evaluates the extent to which each targeted competency has been achieved, supported by evidence from the project lifecycle.

\textbf{CES2.1: Define and manage the requirements of a software system. [In-depth]}

This competency was achieved at an in-depth level. A structured process was conducted to gather and analyze system requirements in collaboration with the Product Owner, resulting in a comprehensive set of functional and non-functional requirements (Section \ref{sec:requirements}). These requirements were continuously prioritized using the MoSCoW methodology, which proved essential for scope management and for delivering a high-value product despite time constraints.

\textbf{CES1.1: Develop, maintain, and evaluate complex and/or critical software systems and services. [Substantially]}

This competency was substantially achieved through the end-to-end development of the full-stack application. The design and implementation of an asynchronous, event-driven processing pipeline (Section \ref{sec:physical-architecture}, Physical Architecture) exemplify the development of a complex system. Additionally, the adoption of Clean Architecture and SOLID principles (Section \ref{sec:logical-architecture}, Logical Architecture) supports maintainability and long-term system evolution.

\textbf{CES1.7: Control quality and design tests in software production. [Substantially]}

Quality assurance was maintained throughout the project lifecycle. Static code analysis tools were initially integrated to monitor code quality (Section \ref{sec:static-code-analysis}). Subsequently, unit and integration tests (Section \ref{sec:unit-and-integration-testing}) were implemented to validate functionality and component interactions. Test coverage was quantitatively assessed, achieving 78\% code coverage. Finally, manual acceptance testing ensured compliance with defined requirements, reflecting a comprehensive approach to quality control.

\textbf{CES1.3: Identify, evaluate, and manage potential risks associated with software construction. [Substantially]}

This competency was substantially achieved, as demonstrated in the Risk Management section (Section \ref{sec:risk-management}). Key risks, including tight deadlines and limited experience with certain technologies, were identified early. Mitigation strategies, such as incorporating time buffers and allocating incidental costs within the project budget (Section \ref{sec:budget-other-costs}), were defined at the project outset. The application of MoSCoW prioritization further supported effective risk management, enabling successful project completion.

\textbf{CES1.9: Demonstrate an understanding of the management and governance of software systems. [Substantially]}

This competency is evidenced by comprehensive project planning and monitoring activities. These include the development of a work breakdown structure (Table \ref{table:summary-of-tasks-with-effort-dependencies-and-resources}), time planning through a Gantt chart (Figure \ref{fig:gantt-diagram}), and detailed budget estimation (Section \ref{sec:budget}). Governance was ensured through regular review meetings with the Product Owner and Thesis Director, maintaining alignment with project objectives. The comparative analysis carried out on the Project Management Results (Section \ref{sec:project-management-results}) between planned and actual outcomes further reflects a mature understanding of project management.

\textbf{CES2.2: Design appropriate solutions in one or more application domains, using software engineering methods that integrate ethical, social, legal, and economic aspects. [Partially]}

This competency was partially achieved. The sustainability analysis (Section \ref{sec:sustainability}) evaluates the project's economic, environmental, and social dimensions, demonstrating consideration of these factors in the solution design. For instance, economic viability (ROI) and social impact (reduction of manual effort) were key drivers. However, certain aspects, particularly deeper legal and ethical considerations such as comprehensive data retention policies, were beyond the project scope, resulting in partial fulfillment of this competency.

